\section{Datamodeller}
\label{sec:datamodeller}

Datamodellene i VideoAPI er utformet for å gi klienten et stabilt og
backend-uavhengig grensesnitt. Siden AntMedia og NHN representerer videorom på
ulike måter, eksponeres ikke råresponsene fra backendene direkte. I stedet mapper
providerne backend-spesifikke svar til felles modeller som kan brukes på tvers av
videomotorene.

Den viktigste responsmodellen er \texttt{VideoRoom}. Modellen representerer et
videorom slik VideoAPI eksponerer det til klienten, med blant annet rom-ID,
provider, status og relevante tilkoblingsfelter. For AntMedia bygges modellen fra
broadcast-data, mens den for NHN bygges fra booking-data. Dette gjør at klienten
kan forholde seg til samme responsstruktur, selv om backendene har ulike
domenemodeller.

Ved opprettelse av rom brukes \texttt{CreateRoomRequest}. Denne modellen
inneholder valgt provider og valgfritt romnavn. Dersom klienten ikke angir navn,
genererer mellomvaren et unikt navn. Dette er i tråd med kravspesifikasjonen,
som beskriver at mellomvaren skal sette navn eller ID dersom dette ikke sendes
med av klienten~\cite{hdo_kravspec}.

For tilgang til videorom brukes \texttt{WebRtcLinkResponse}. Denne modellen
normaliserer to ulike tilgangsformer. AntMedia returnerer en kortlivet
tokenbasert tilgang til en stream, mens NHN returnerer en møte-URL med nødvendig
tilgangsinformasjon. Modellen gir derfor klienten ett felles svarformat, men
uten å skjule at verdien må brukes forskjellig avhengig av backend.

Datamodellene er dermed et kompromiss mellom standardisering og presisjon. De
skjuler unødvendige forskjeller fra klienten, men lar backend-spesifikke
forskjeller komme frem der de påvirker hvordan operasjonen faktisk brukes.